\section{Duplicate 与 Intervention：把 learned lookup 拆成可观察步骤}

真实数据上的 corruption 说明关系通道被使用，却仍不清楚模型学到了什么算法。Semi-synthetic duplicate task 构造一个更透明的机制：复制若干 rows，把 originals 的 target 遮住、duplicates 的 target 暴露；正确预测必须匹配对应 row、读取 target、再复制回 query。

\lecturefigure{slide-27-duplicate-intervention.jpg}{Semi-synthetic duplicate task、attention、预测与 target intervention 的完整证据链。}{Stanford Online 官方视频 00:54:44--01:00:41。}

\begin{knowledgebox}{读图：五个 panel 的推理顺序}
\begin{enumerate}
  \item (a) 构造 original/duplicate rows，并遮住 originals 的 targets；
  \item (b) 检查 ABD attention 是否集中到对应 duplicates；
  \item (c) 验证 predictions 与 duplicate targets 高度相关；
  \item (d) 在 test time 主动改变某个 duplicate target，不重新训练；
  \item (e) 观察 original prediction 是否沿 intervention 变化。
\end{enumerate}
\end{knowledgebox}

这个任务要求模型学会一个组合程序：
\[
x_i\xrightarrow{\text{match features}}
x_{j(i)}^{\mathrm{dup}}
\xrightarrow{\text{read target}}
y_{j(i)}^{\mathrm{dup}}
\xrightarrow{\text{copy/transform}}
\widehat y_i.
\]

\begin{itemize}
  \item $j(i)$：与 original row $i$ 对应的 duplicate index；
  \item match 由 learned attention representation 完成，而不是预写 row ID；
  \item intervention 改变 $y_{j(i)}^{\mathrm{dup}}$，其他机制保持不变；
  \item 若 $\widehat y_i$ 随 intervention 稳定移动，说明模型使用了可泛化的 lookup path。
\end{itemize}

\teachervoice{讲者报告 intervention 后 prediction 与 duplicate value 的相关性仍约为 99.6\%。关键是模型未重新训练，却跟随一个训练分布之外的 target value 改变；这比只看普通 test error 更接近机制检验。}

\subsection{为什么它不只是 nearest neighbor}

Plain k-NN 也能找到 duplicate 并复制 target，因此基础任务还不足以展示 NPT 的可学习关系变换。作者进一步把 duplicate target 加一；固定 copy rule 会系统性出错，而 NPT 可以学习先 lookup 再减一。更一般地，NPT 可以学习 joint feature--target relation，而非只在固定 metric 下做局部平均。

\begin{importantbox}{NPT 学到的不是一个静态邻居表}
最强证据不是“attention 看向相似 row”，而是模型在 match 之后还能执行 task-dependent transformation。多层 ABD/ABA 使 lookup、attribute reasoning 与 value transformation 可以组成一段 learned algorithm。
\end{importantbox}

\teachervoice{Kossen 主动承认基础 duplicate task 可被 nearest neighbor 解决，然后提出加一变换来提高难度。这种论证方式值得学习：不要回避简单 baseline，而要构造能区分机制能力的最小实验。}

\begin{warningbox}{Semi-synthetic success 的外推边界}
任务被刻意设计为存在清晰 duplicate 与可验证 intervention，因此适合证明 representational capability。现实数据通常没有完美匹配、干净单一机制或可操控 labels；成功不等价于 NPT 已掌握一般 causal reasoning。
\end{warningbox}

\subsection{本章小结}

Duplicate/intervention experiment 将跨 datapoint reasoning 拆成 match、read、transform、write 四步，并用 test-time intervention 验证 prediction 随 context 改变。它是强机制证据，但仍局限于透明的 semi-synthetic construction。

